\documentclass{article}
\usepackage{amsmath, amssymb, amsthm}
\usepackage{hyperref}
\usepackage{geometry}
\geometry{margin=1in}

\title{Erd\H{o}s Problem \#340}
\date{Last updated: 2025-08-31}
\author{Source: erdosproblems.com}

\begin{document}
\maketitle

\noindent\textbf{Status:} OPEN \quad \textbf{No prize}

\noindent\textbf{Tags:} number theory, additive combinatorics, sidon sets

\medskip

\noindent\textbf{Problem Statement:}

Let $A=\{1,2,4,8,13,21,31,45,66,81,97,\ldots\}$ be the greedy Sidon sequence: we begin with $1$ and iteratively include the next smallest integer that preserves the Sidon property (i.e. there are no non-trivial solutions to $a+b=c+d$). What is the order of growth of $A$? Is it true that\[\lvert A\cap \{1,\ldots,N\}\rvert \gg N^{1/2-\epsilon}\]for all $\epsilon&#62;0$ and large $N$?




This sequence is sometimes called the Mian-Chowla sequence. It is trivial that this sequence grows at least like $\gg N^{1/3}$. Erd\H{o}s and Graham \cite{ErGr80} also asked about the difference set $A-A$, whether this has positive density, and whether this contains $22$. It does contain $22$, since $a_{15}-a_{14}=204-182=22$. The smallest integer which is unknown to be in $A-A$ is $33$ (see A080200 ). It may be true that all or almost all integers are in $A-A$. This sequence is at OEIS A005282 . See also [156] .




References

[ErGr80] Erd\H{o}s, P. and Graham, R., Old and new problems and results in combinatorial number theory . Monographies de L'Enseignement Mathematique (1980).

\noindent\textbf{Related OEIS sequences:} A080200, A005282


\bigskip
\noindent\small{Source: \url{https://www.erdosproblems.com/340}}
\end{document}
